
\section{Eficiência energética}

Com o crescimento da demanda por serviços em nuvem e processamento distribuído, o consumo e eficiência de Data Centers se mostrou como um desafio de engenharia. Para quantificar e padronizar a eficiência da infraestrutura física que abriga esses sistemas, o consórcio The Green Grid e a norma internacional ISO/IEC 30134-2 estabeleceram a métrica Power Usage Effectiveness (PUE) como o principal indicador de desempenho energético do setor \cite{iso30134pue}. O PUE é definido pela razão entre a energia total consumida pela instalação e a energia efetivamente entregue aos equipamentos de Tecnologia da Informação:

\begin{center}
	\footnotesize
	\begin{equation}
		\text{PUE} = \frac{P_{\text{Total da Instalação}}}{P_{\text{Equipamentos de TI}}} = \frac{P_{\text{TI}} + P_{\text{Refrigeração}} + P_{\text{Perdas Elétricas}} + P_{\text{Iluminação/Auxiliares}}}{P_{\text{TI}}}
	\end{equation}  
\end{center}

Em um cenário ideal, em que toda a energia fornecida ao edifício fosse integralmente consumida por servidores, ativos de armazenamento e comutadores de rede, o índice PUE seria igual a 1,0.
Embora instalações de grande escala alcancem valores próximos a 1,1 ou 1,2 (ou seja, a instalação completa custa mais que somente os equipamentos de TI), a métrica apresenta uma limitação: o PUE avalia exclusivamente a eficiência da infraestrutura de suporte predial (climatização e conversão elétrica), mas não mensura a eficiência computacional interna dos servidores \cite{dayarathna2016data}.

Essa distinção torna-se evidente ao analisar o fenômeno da não proporcionalidade energética do hardware de servidor. Servidores físicos operando em estado ocioso ou sob baixas taxas de utilização de CPU (entre 10\% e 20\%) ainda consomem cerca de 50\% a 60\% de sua potência elétrica máxima nominal devido à energização contínua de barramentos, memórias e fontes de alimentação \cite{beloglazov2012energy}. Para eliminar esse desperdício de potência ociosa, Data Centers empregam algoritmos de consolidação dinâmica de servidores. Utilizando heurísticas de empacotamento aplicadas a máquinas virtuais e orquestradores de contêineres, as cargas de trabalho dispersas são migradas e concentradas no menor número possível de nós físicos ativos, permitindo que os servidores excedentes sejam colocados em estados de suspensão ou desligados, utilizando o recurso da elasticidade segundo a demanda. Porém, abaixar o denominador da equação pode acabar aumentando o valor total do PUE se a climatização e auxiliares não forem modulados junto com os servidores.

Consequentemente, a gestão térmica atua como o pilar complementar indispensável da eficiência energética, visto que os sistemas de refrigeração representam historicamente a maior parcela de
consumo não-TI de um Data Center. A consolidação de cargas via contêineres eleva a utilização média das CPUs e intensifica a densidade de dissipação térmica por gabinete, exigindo arquiteturas de rejeição de calor alinhadas às diretrizes ambientais do comitê técnico ASHRAE TC 9.9 \cite{ashrae2021thermal}. A otimização térmica moderna estrutura-se em três pontos: (i) o isolamento aerodinâmico por meio da contenção de corredores frios ou quentes, que impede a recirculação e a mistura entre o ar insuflado e o ar de exaustão, permitindo elevar os pontos de ajuste de temperatura de entrada sem risco de formação de pontos quentes; (ii) a adoção
de sistemas de \emph{Free Cooling} (economizadores a ar ou a água), que aproveitam condições climáticas externas favoráveis para resfriar o fluido térmico sem o acionamento contínuo de compressores mecânicos; e (iii) o resfriamento líquido direto no componente (\emph{Direct-to-Chip Liquid Cooling}) ou por imersão, cuja capacidade térmica volumétrica superior à do ar viabiliza a operação eficiente de racks de alta densidade \cite{ashrae2021thermal}. A Tabela \ref{tab:sintese} sintetiza a correlação entre os mecanismos de segurança em contêineres e a eficiência energética da instalação.

\begin{table}[h!]
	\centering
	\caption{Síntese integradora entre segurança de contêineres e eficiência energética em Data Centers.}
	\label{tab:sintese}
	\footnotesize
	\renewcommand{\arraystretch}{1.2}
	\begin{tabularx}{\textwidth}{@{} l X X X @{}}
		\toprule
		\textbf{Dimensão} & \textbf{Desafio Técnico} & \textbf{Mecanismo / Solução} & \textbf{Impacto Cruzado (Segurança $\times$ Energia)} \\
		\midrule
		Isolamento de Kernel & Superfície de ataque e escalada de privilégios & \textit{Namespaces}, \textit{cgroups}, \textit{Seccomp} e \textit{Rootless} & Viabiliza alta densidade de consolidação sem o gasto energético de VMs completas. \\
		Imagens e Rede & CVEs em camadas base e tráfego \textit{East-West} & Imagens \textit{Distroless}, \textit{Network Policies} e \textit{eBPF} & Binários enxutos e filtragem no \textit{kernel} via \textit{eBPF} reduzem ciclos de CPU e rede. \\
		Alocação de Nós & Desperdício por potência ociosa (\textit{Idle Power}) & Consolidação dinâmica (\textit{Bin Packing}) de cargas & Reduz a potência total de TI ($P_{\text{TI}}$), exigindo contenção estrita de \textit{multitenancy}. \\
		Gestão Térmica & Dissipação de calor em \textit{racks} densos e PUE alto & Contenção de corredores, \textit{Free} e \textit{Liquid Cooling} & Reduz a potência de refrigeração no numerador do PUE sob alta consolidação. \\
		\bottomrule
	\end{tabularx}
\end{table}

A urgência em otimizar o PUE, consolidar servidores e modernizar a gestão térmica é amplificada pela recente expansão global de \emph{data centers} de hiperescala, impulsionada pela computação em nuvem e por cargas intensivas de inteligência artificial. Relatórios da Agência Internacional de Energia (IEA) projetam que o consumo elétrico global do setor atingirá 945~TWh até 2030, um volume equivalente à demanda total de eletricidade do Japão, enquanto em polos consolidados, como a Irlanda, essas instalações já absorveram 23\% de toda a eletricidade medida no país em 2025 \cite{iea2025energy, cso2025ireland}. Nos Estados Unidos, levantamentos do \emph{Lawrence Berkeley National Laboratory} (LBNL) apontam um consumo de 192 TWh em 2024 (4,7\% da matriz elétrica nacional), com projeção de ascensão para 649 TWh até 2030 (11,8\% da demanda nacional); contudo, observa-se uma diferença operacional: enquanto a média nacional geral registra um PUE de 1,45 (projetado para 1,36 em 2030), instalações dedicadas a cargas de IA impulsionadas pelo resfriamento líquido alcançam um PUE médio estimado de 1,145 (caindo para 1,136 até 2030) \cite{shehabi2026lbnl}. No cenário brasileiro, essa pressão energética reflete-se nas projeções do Operador Nacional do Sistema Elétrico (ONS), CCEE e EPE, que contabilizam 27 contratos assinados e 30 pareceres favoráveis de acesso à Rede Básica até julho de 2026 (além de 21 em análise), projetando que a carga exclusiva de data centers no Sistema Interligado Nacional (SIN) salte de 315 MW médios em 2026 para 5.653 MW médios até 2030  \cite{ons2026datacenters}. Como essas novas instalações operam com densidades térmicas que saltaram de 10~kW para mais de 40 a 100~kW por \emph{rack}, o resfriamento a ar tradicional torna-se termodinamicamente insuficiente, tornando obrigatória a combinação de refrigeração líquida direta ao chip com a consolidação agressiva de contêineres para mitigar o desperdício em megawatts de potência ociosa.
